\chapter{Chemotherapy}

\section*{Overview}
Chemotherapy is the use of drugs to destroy cancer cells. It is one of the main types of cancer treatment, along with surgery and radiation therapy. Chemotherapy works by targeting rapidly dividing cells, which is a characteristic of cancer cells. The drugs travel through the bloodstream to reach cancer cells throughout the body, making it effective for cancers that have spread or are likely to spread.

\section*{Symptoms}
Chemotherapy side effects occur because the drugs also affect healthy cells that divide rapidly. Common side effects include:
\begin{itemize}
    \item Fatigue and weakness
    \item Nausea and vomiting
    \item Hair loss (alopecia)
    \item Mouth sores
    \item Loss of appetite
    \item Increased risk of infection (due to low white blood cell count)
    \item Easy bruising or bleeding (due to low platelet count)
    \item Anemia (fatigue, shortness of breath)
    \item Diarrhea or constipation
    \item Skin and nail changes
    \item "Chemo brain" (cognitive changes)
\end{itemize}
Side effects vary depending on the drugs used, dosage, and individual patient factors.


\section*{Causes}
Chemotherapy is prescribed by oncologists to:
\begin{itemize}
    \item \textbf{Cure cancer:} Eliminate all cancer cells
    \item \textbf{Control cancer:} Slow or stop growth of advanced cancer
    \item \textbf{Relieve symptoms:} Palliative care to improve quality of life
    \item \textbf{Prepare for surgery:} Shrink tumors before operation (neoadjuvant)
    \item \textbf{Prevent recurrence:} Kill remaining cells after surgery (adjuvant)
\end{itemize}
The choice of chemotherapy drugs depends on cancer type, stage, genetic markers, and patient health status.


\section*{Diagnosis}
Before starting chemotherapy, oncologists determine:
\begin{itemize}
    \item Cancer type and stage through biopsy and imaging
    \item Whether chemotherapy is appropriate for the specific cancer
    \item Which drug regimen will be most effective
    \item Patient's ability to tolerate treatment (organ function tests)
    \item Potential drug interactions with current medications
\end{itemize}
Genetic testing of tumors helps identify targeted therapies and predict response.


\section*{Treatment}
Chemotherapy can be given in several ways:
\begin{itemize}
    \item \textbf{Intravenous (IV):} Most common, through a vein
    \item \textbf{Oral:} Pills or liquids taken by mouth
    \item \textbf{Intramuscular/Subcutaneous:} Injection into muscle or under skin
    \item \textbf{Intrathecal:} Directly into spinal fluid
    \item \textbf{Topical:} Cream applied to skin
    \item \textbf{Intracavitary:} Into body cavities (abdomen, bladder)
\end{itemize}
Treatment is given in cycles with rest periods to allow healthy cells to recover. Common regimens include combinations of drugs like cyclophosphamide, doxorubicin, cisplatin, and paclitaxel.


\section*{Prevention}
While chemotherapy cannot be prevented, patients can:
\begin{itemize}
    \item Discuss all medications with oncologist to avoid interactions
    \item Stay hydrated before and after treatment
    \item Maintain good nutrition to support recovery
    \item Practice good hygiene to prevent infections
    \item Report side effects promptly to healthcare team
    \item Use cold caps to reduce hair loss (where appropriate)
    \item Avoid pregnancy during treatment (teratogenic risk)
\end{itemize}
Proper preparation and supportive care minimize complications.


\clearpage